% ACL 2023 format: 5-page body, then Limitations, Ethics and references.
% Synthesises talk_like_a_graph.v3 (procedure, graph-blind solver) and
% compare/independent (target-verified rescoring, C/J, RWSE resolution).
% Every number is taken from one of those two sources: ../NUMBERS.md (v3,
% terminated pairs) and ../compare/independent/CLAIM_INDEX.md (S, capped =
% failure).
\documentclass[11pt]{article}

\usepackage{acl2023}

\usepackage{times}
\usepackage{latexsym}
\usepackage[T1]{fontenc}
\usepackage[utf8]{inputenc}
\usepackage{microtype}
\usepackage{inconsolata}
\usepackage{graphicx}
\usepackage{booktabs}
\usepackage{amsmath}

\newcommand{\task}[1]{\texttt{#1}}
\newcommand{\cond}[1]{\texttt{#1}}

\title{Structural Primers Trigger Procedural Shortcuts in LLMs:\\
Answer Leakage, Invariants and Feature Resolution on Graphs}

\author{
  Gal Barak \\
  Tel Aviv University \\
  \footnotesize\texttt{github.com/GalBrk} \\\And
  Nitzan Zacharia \\
  Tel Aviv University \\
  \footnotesize\texttt{github.com/NitzanZacharia} \\\And
  Inbal Moryles \\
  Tel Aviv University \\
  \footnotesize\texttt{github.com/iinbal} \\
}

\begin{document}
\maketitle
\raggedbottom

\begin{abstract}
Graph neural networks routinely receive structural encodings such as node
degree, clustering and random-walk return probabilities as input features.
We serialize the same descriptors as a textual \emph{primer} before a graph
given to a language model. Does the primer change how the model executes a
graph algorithm, and which structural distinctions survive serialization? We evaluate Qwen3-1.7B and Qwen3-4B,
with plain and thinking modes, on 40-node Erd\H{o}s--R\'enyi graphs at seven
edge densities ($p{=}.10$ to $.85$), over 84,000 generations with gold
answers verified against reconstructed graphs. Two controls separate access to an answer from its use: a graph-blind
solver flags primers that leak the answer, and a paired diagnostic checks
a model's degree and neighborhood answers for one node against each other. A degree primer lifts plain
Qwen3-4B's edge counting from 2.00\% to 29.25\%, yet lowers its joint
degree--neighborhood correctness from 85.75\% to 76.50\%. On degree queries
the same model skips the neighbor list in 44--74\% of responses. It gains
27 points at $p{=}.85$, where its own count is 33\% accurate, and loses up
to 12 where the count is 99\% accurate. Qwen3-1.7B with thinking instead
keeps counting and uses the primer as a check (+16.0 points). Printed to two
decimals, random-walk features take on average 3.30 distinct values across a
graph's 40 nodes at $p{=}.50$, fewer than degree alone. A leaked answer can
change which procedure a model runs, and its effect follows the accuracy of
the procedure it displaces. Neither route is evidence of graph
understanding.
\end{abstract}

\section{Introduction}

A language model (LLM) reads a graph given as text token by token. To
answer ``what is the degree of node 16?'' it must find node 16's adjacency
line and count it. \citet{fatemi2023talklikeagraph} showed that accuracy on
such GraphQA tasks depends strongly on the textual encoding. Graph learning
has a standard remedy for a model that cannot easily compute a structural
quantity: supply it. Degree, clustering and random-walk structural encodings
(RWSE) are input features of graph neural networks (GNNs) and graph
transformers \citep{dwivedi2022graph,rampasek2022recipe}. We study the
textual analogue, a \emph{structural primer}: one sentence per node, placed
before the graph, stating its degree, clustering coefficient or RWSE.

\paragraph{Research question.} \emph{How do textual structural primers alter
an LLM's procedural execution of graph tasks, and what are the expressivity
limits of these primers in preserving structural invariants?} Aggregate
accuracy cannot answer it: it conflates \emph{access} to a
text feature that already states the answer, \emph{execution} of the
algorithm the task requires, and \emph{consistency} with relations any
correct answer must satisfy. A primer that states every degree turns
\task{node\_degree} into lookup, so a gain there says nothing about whether
the model counts. A model that looks up one answer and counts another can
violate the identity $d_v{=}|N(v)|$, which holds on every simple graph. Nor
does accuracy show whether a printed feature still separates the nodes its
exact values separate, which bounds what a primer can convey at all.

\paragraph{Contributions.} (1)~A \emph{graph-blind solver} that reads only the
rendered primer and node/edge counts and marks where a primer leaks the answer
(\S\ref{sec:solver}). (2)~A \emph{paired consistency diagnostic} that scores
a model's degree and neighborhood answers for the same node against
$d_v{=}|N(v)|$, separating agreement ($C$) from joint correctness ($J$)
(\S\ref{sec:cj}). (3)~A classification of responses into retrieval and
enumeration, which shows the primer replacing or checking the counting
procedure (\S\ref{sec:procedure}). (4)~A count of the distinct values a
two-decimal RWSE primer takes per graph, which from $p{=}.35$ falls below
that of degree (\S\ref{sec:resolution}).

\section{Related Work}

\paragraph{Graphs as text.} GraphQA \citep{fatemi2023talklikeagraph}, NLGraph
\citep{wang2023nlgraph}, GPT4Graph \citep{guo2023gpt4graph} and GraphArena
\citep{tang2025grapharena} vary encodings, prompts and tasks; NLGift
\citep{zhang2024nlgift} asks whether gains from tuning on synthetic graph
data generalize. \citet{sanford2024understanding} show that a single-layer
transformer can represent node degree and edge count, suggesting that
difficulty lies in execution, not architectural capacity.
\citet{huang2024can} find that neighborhood text helps node classification
through label-relevant phrases rather than structure.
\citet{dziri2023faithandfate} show that transformers solve compositional
tasks through shortcut subgraphs of the computation. Graph-QA evaluations
mostly score final answers. We know of none that records the procedure
behind an answer or checks two answers about one node against each other.

\paragraph{Structural encodings.} Message-passing GNNs are bounded by the
1-WL test \citep{xu2019powerful}, whose first refinement round assigns each
node its degree. RWSE and related encodings add information beyond that
bound \citep{dwivedi2022graph,rampasek2022recipe}. A GNN receives them as
real-valued vectors, in which small differences between nodes of equal
degree stay distinct. GraphToken \citep{perozzi2024letyourgraph} instead
learns a graph encoder for a frozen LLM. A textual primer needs no training,
but it passes every value through the tokenizer as a fixed-precision digit
string. We know of no measurement of how much node-level resolution this
serialization keeps.

\paragraph{Shortcuts and partial inputs.} Shortcut learning
\citep{geirhos2020shortcut} and hypothesis-only baselines
\citep{gururangan2018annotation,poliak2018hypothesis} show that high accuracy
can come from cues that do not require the intended computation. A partial
input's failure does not rule a shortcut out \citep{feng2019misleading}. We
apply the partial-input idea to an intervention (the primer) rather than a
dataset. Added text can itself degrade reasoning
\citep{shi2023distracted,levy2024sametask}, so we include a content-free
length control.

\section{Method}

\subsection{Corpus and models}
\label{sec:corpus}

Graphs are Erd\H{o}s--R\'enyi $G(n,p)$
\citep{erdosrenyi1959random,gilbert1959random} with $n{=}40$, 100
independently seeded graphs per density. The main sweep covers
$p\in\{.10,.20,.35,.50\}$, six GraphQA tasks and seven conditions (67,200
generations). An extension adds $p\in\{.65,.75,.85\}$ for \task{node\_degree}
and \task{edge\_existence} only (16,800). \task{node\_degree},
\task{connected\_nodes} and \task{edge\_existence} query the same node of a
graph. Four arms answer every prompt: Qwen3-1.7B and Qwen3-4B
\citep{qwen2025qwen3} with thinking off (1.7B, 4B) and on (1.7B-T, 4B-T).
Decoding is greedy with an 8,192-token budget (2,048 for the plain arms in the
extension). Every saved gold answer matches the graph rebuilt from its
seed; the 84,000 unique generations form 840 cells of 100. Every graph here has 40 nodes and a cycle, so \task{node\_count} and
\task{cycle\_check} have a constant gold and are excluded from comparisons.

\subsection{Primers}

A prompt consists of the primer, the graph in the \texttt{incident} encoding
of \citet{fatemi2023talklikeagraph}, and the question. All seven conditions
share one renderer (Table~\ref{tab:conditions}). Clustering and RWSE, the
return probabilities $(P^2_{vv},P^3_{vv})$ of the random walk $P{=}D^{-1}A$
(0 for isolated nodes), are printed to two decimals.
\cond{filler} restates the node set (``Node $i$ is simply present within the
graph $G$''), which the encoding already lists, at 12\% above the length of
\cond{clustering}.

\begin{table}[t]
\centering
\small
\setlength{\tabcolsep}{2.4pt}
\begin{tabular}{@{}llrrr@{}}
\toprule
Condition & States & Chars & deg. & e.cnt \\
\midrule
\cond{none}       & no primer                   & --    & 14.4 & 3.0 \\
\cond{filler}     & ``Node $i$ is simply present'' & 1,829 & 14.4 & 3.0 \\
\cond{components} & component count (graph)     & 37    & 14.4 & 3.0 \\
\cond{clustering} & clustering coefficient      & 1,629 & 14.4 & 10.4 \\
\cond{rwse}       & return prob., 2 and 3 steps & 2,949 & 21.0 & 3.0 \\
\cond{degree}     & degree                      & 891   & \textbf{100} & \textbf{100} \\
\cond{all}        & the three rows above        & 4,691 & \textbf{100} & \textbf{100} \\
\bottomrule
\end{tabular}
\caption{Conditions (per node unless marked), mean length, and graph-blind
solver accuracy (\%) on \task{node\_degree} (deg.) and \task{edge\_count} (e.cnt),
over seven densities, 300 held-out graphs each. \textbf{Bold}:
answer-carrying. No condition exceeds 0.81\% on \task{connected\_nodes} or
74.6\% on \task{edge\_existence}.}
\label{tab:conditions}
\end{table}

\subsection{Graph-blind solver}
\label{sec:solver}

The solver reads the rendered primer text and the graph's node and edge
counts (except the count a task asks for), never the edge list. It applies
16 exact rules where they exist, and otherwise one heuristic or 8 rules
fitted on graphs disjoint from those scored. Its accuracy is therefore a
lower bound on what the primer and counts determine. A primer is
\emph{answer-carrying} on a task when the solver scores at least 85\% there,
and \emph{side information} otherwise. Any threshold in $(74.6,100]$ gives
the same split. Only \cond{degree} and \cond{all} carry answers, on
\task{node\_degree} (the queried degree) and \task{edge\_count}
($|E|{=}\tfrac12\sum_v d_v$).

\subsection{Scoring and inference}

The primary outcome is $S{=}\mathbf{1}[\hat y{=}y]\cdot\mathbf{1}[\text{not
capped}]$: exact match (set equality for neighbors), with a generation that
reaches its budget scored as a failure. Main-sweep effects pool the four
densities ($n{=}400$ paired graphs). Other effects and all response shares
(\S\ref{sec:procedure}, \S\ref{sec:side}) use terminated generations,
dropping a pair when either generation is capped; this coincides with $S$
wherever nothing is capped, as for plain 4B on \task{node\_degree}. Effects
are paired differences in percentage points (pp) against \cond{none}, with
exact McNemar tests \citep{mcnemar1947note} and Benjamini--Hochberg
adjustment \citep{benjaminihochberg1995controlling} over the six primer
contrasts per (arm, task), or (arm, metric) for $C$ and $J$. Intervals are
95\% paired bootstrap intervals over graphs, stratified by density
\citep{efron1979bootstrap}. All tests are exploratory; intervals are not
simultaneous.

\subsection{Cross-query consistency}
\label{sec:cj}

For node $v$ let $\hat d_v$ and $\hat N_v$ be the model's answers to the
separate \task{node\_degree} and \task{connected\_nodes} prompts. $C_v{=}1$
when both complete, $\hat d_v$ parses as an integer, the predicted neighbors
are valid node ids, and $\hat d_v{=}|\hat N_v|$. In that case the two answers
satisfy the degree identity. $J_v{=}S_{v,\text{deg}}\,S_{v,\text{nbr}}$ is
joint correctness. Since $J_v\le C_v$, $C{-}J$ is the rate of
\emph{consistent but wrong} pairs. $C$ tests a necessary relational
invariant at the output level, not an internal graph representation.

\subsection{Procedure classification}
\label{sec:procedure-def}

An \emph{enumerating} \task{node\_degree} response restates the queried node's neighbor list and
itemizes at least five entries, one per line. A \emph{retrieving} response
answers without restating the list. A retrieval need not return the stated
degree; accuracy among retrievals measures how often it does. A thinking
response \emph{reports a discrepancy} when it uses a discrepancy word
(discrepancy, contradiction, conflict, inconsistency) outside a
hypothetical. Without a primer there is nothing to retrieve, so accuracy
under \cond{none} is the accuracy of the model's own count.

\begin{figure*}[t]
\centering
\includegraphics[width=\textwidth]{density_diagnostics.pdf}
\caption{Density diagnostics (100 graphs per point; shading: the $p\ge.65$
extension, plain-arm budget 2,048 tokens). (a)~Plain 4B \task{node\_degree}
success, \cond{none} and \cond{degree}. (b)~Plain 1.7B
\task{edge\_existence}, no primer: raw, balanced and majority-label success.
(c)~Mean distinct RWSE pairs among a graph's 40 nodes at 2, 4 or 6
decimals; models saw two decimals.}
\label{fig:density}
\end{figure*}

\section{Results}

\begin{table}[t]
\centering
\footnotesize
\setlength{\tabcolsep}{1.5pt}
\resizebox{\columnwidth}{!}{%
\begin{tabular}{@{}llrrrrrrr@{}}
\toprule
Arm & Task & None & Fill. & Comp. & Clust. & RWSE & Deg. & All \\
\midrule
1.7B & Deg & 60.50 & $-0.50$ & 1.75 & 3.00 & 4.00 & 7.50$^{*}$ & 10.25$^{*}$ \\
 & Nbr & 72.50 & $-6.25^{*}$ & 0.75 & 1.25 & 1.00 & 3.50 & 2.00 \\
 & EE & 70.50 & $-15.75^{*}$ & $-8.25^{*}$ & 3.75 & 5.00 & 4.00 & 11.50$^{*}$ \\
 & EC & 1.25 & $-0.50$ & $-1.00$ & $-0.50$ & $-0.75$ & 3.00 & 0.75 \\
\addlinespace[2pt]
1.7B-T & Deg & 76.25 & $-3.00$ & $-1.00$ & 1.50 & 2.50 & 8.75$^{*}$ & 12.25$^{*}$ \\
 & Nbr & 78.75 & 2.50 & 2.50 & 2.00 & 0.25 & 6.25$^{*}$ & 6.25$^{*}$ \\
 & EE & 99.00 & $-0.50$ & $-0.75$ & $-1.00$ & $-1.25$ & $-2.50$ & $-1.75$ \\
 & EC & 15.75 & $-1.50$ & $-0.75$ & $-1.25$ & $-3.25$ & $-3.25$ & 5.00 \\
\addlinespace[2pt]
4B & Deg & 99.25 & $-0.50$ & 0.25 & $-0.50$ & $-2.25^{*}$ & $-6.50^{*}$ & $-8.00^{*}$ \\
 & Nbr & 86.00 & $-12.50^{*}$ & 9.25$^{*}$ & 1.00 & 0.00 & $-3.00$ & $-2.75$ \\
 & EE & 90.50 & $-7.00^{*}$ & $-2.75$ & 4.00$^{*}$ & $-10.00^{*}$ & $-4.75^{*}$ & $-10.00^{*}$ \\
 & EC & 2.00 & 1.25 & $-1.25$ & $-1.25$ & $-1.50$ & 27.25$^{*}$ & 3.75$^{*}$ \\
\addlinespace[2pt]
4B-T & Deg & 97.00 & 2.00 & 0.25 & 1.75 & $-1.25$ & 2.50 & 0.50 \\
 & Nbr & 95.50 & 1.25 & 0.50 & $-0.75$ & 0.50 & 0.50 & $-1.25$ \\
 & EE & 99.75 & 0.00 & $-0.50$ & $-0.75$ & $-1.25$ & 0.00 & $-1.00$ \\
 & EC & 20.00 & 12.25$^{*}$ & 24.75$^{*}$ & 24.75$^{*}$ & 10.25$^{*}$ & 8.50$^{*}$ & 1.75 \\
\bottomrule
\end{tabular}}
\caption{Main sweep ($p\le.50$, $n{=}400$ per cell): success $S$ (\%;
capped = 0) under \cond{none} and paired effects (pp) on
\task{node\_degree} (Deg), \task{connected\_nodes} (Nbr),
\task{edge\_existence} (EE) and \task{edge\_count} (EC). $^{*}$: BH-adjusted $q<.05$.}
\label{tab:main}
\end{table}

\subsection{Answer access is not answer use}

Table~\ref{tab:main} shows that an answer-carrying primer does not
reliably raise accuracy on the task it leaks. \cond{degree} reduces
\task{edge\_count} to a 40-term sum. Plain 4B sets out to sum degrees in
96--100\% of responses even without a primer. With \cond{degree} it goes
from 2.00\% to 29.25\% (+27.25 [22.75, 31.75]; 114 graphs fixed, 5 broken).
Plain 1.7B gains only +3.00 ($q{=}.101$). Its sums move closer (mean
absolute error 136 to 40 edges, terminated generations) but are rarely
exact. 1.7B-T changes by $-3.25$. The thinking arms rarely complete an edge count. Without a primer 4B-T
finishes 20.50\% of generations and is correct on 97.56\% of those, for
20.00\% success. Under \cond{degree} it completes 41.25\% but succeeds on
only 28.50\%. \cond{components} and \cond{clustering} carry no edge count
(solver 3.0\% and 10.4\%), yet both reach 44.75\%, 12.50 pp above
\cond{filler} (32.25\%; intervals [6.75, 18.01] and [7.50, 17.75]).
\cond{components} states one component for every graph at $p\ge.20$, and
there it still beats \cond{filler} (64/35/8\% against 46/22/1\% at
$p{=}.20/.35/.50$). Its gain acts mainly on whether the thinking
arm completes (45.50\% against 33.50\% under \cond{filler}), not on
graph-specific content.

\subsection{Primers replace or check the counting procedure}
\label{sec:procedure}

Table~\ref{tab:procedure} follows plain 4B on \task{node\_degree} across all
seven densities. Without a primer its count is 99--100\% accurate up to
$p{=}.50$ and falls to 33\% at $.85$. Under
\cond{degree} it often skips the neighbor list, retrieving in 44--74\% of responses
at 73--98\% accuracy. At $p{=}.35$--$.75$, enumeration falls from
85--100\% of responses to at most 8\%. The primer's effect follows the accuracy of
the count it replaces: $-6$, $-1$, $-7$, $-12$ and $-6$ pp at
$p{=}.10$--$.65$, then $+7$ and $+27$ at $.75$ and $.85$ (100 graphs each;
Figure~\ref{fig:density}a). At $p{=}.50$ the primer turns 13 correct answers
into errors and corrects one. \cond{filler}, twice as long, costs 2 pp, so
the loss is not a cost of added text. From $p{=}.50$ the responses that do
not retrieve are only 79\%, 45\%, 28\% and 16\% accurate: 68--92\% of them
open with an answer, usually not the stated degree, and then recount.
Under \cond{all}, which states the degree among two other statistics,
retrieval falls to 15--51\% and accuracy drops at every density, by up to 37 pp at
$p{=}.65$.

The same primer induces a different procedure in 1.7B-T. It retrieves in at
most 2\% of responses and, from $p{=}.35$, enumerates in 76--98\%. It uses a
discrepancy word in 23\% of responses up to $p{=}.50$ and 49\% beyond (0.3\%
and 4\% without a primer), and is correct on 61\% and 69\% of those. Over the
seven densities \cond{degree} adds +16.0 pp [12.0, 19.7] with capped
generations scored as errors (+19.9 [16.3, 23.6] over terminated pairs). It
also raises capping from 1.1\% to 7.6\%. Plain 1.7B retrieves in at most 4\%
of responses. It gains +10 and +18 pp only at $p{=}.20$ and $.35$, where its
own count is 80\% and 41\% accurate. From $p{=}.50$, with baselines of
5--30\%, it stays within 2 pp. As a \emph{replacement}, the leaked answer
helps or hurts according to the accuracy of the count it displaces. As a
\emph{check} on a count 1.7B-T still executes, it helps.

\begin{table}[t]
\centering
\small
\setlength{\tabcolsep}{2.6pt}
\begin{tabular}{lrrrrrrr}
\toprule
$p$ & .10 & .20 & .35 & .50 & .65 & .75 & .85 \\
\midrule
acc., \cond{none}      & 100 & 99 & 99 & 99 & 82 & 50 & 33 \\
\quad enumerates       & 1 & 34 & 95 & 100 & 94 & 85 & 44 \\
acc., \cond{degree}    & 94 & 98 & 92 & 87 & 76 & 57 & 60 \\
\quad retrieves        & 44 & 56 & 74 & 62 & 71 & 64 & 69 \\
\quad acc.\ retrievals & 89 & 98 & 91 & 92 & 89 & 73 & 80 \\
\quad acc.\ the rest   & 98 & 98 & 96 & 79 & 45 & 28 & 16 \\
\quad enumerates       & 2 & 16 & 6 & 8 & 2 & 1 & 0 \\
\bottomrule
\end{tabular}
\caption{Plain Qwen3-4B on \task{node\_degree} by density (\%, 100 graphs
each, none capped). Procedures as in \S\ref{sec:procedure-def}; the five-entry
rule undercounts enumeration at $p\le.20$. \emph{The rest}: responses that
do not retrieve.}
\label{tab:procedure}
\end{table}

\subsection{Agreement does not certify correctness}
\label{sec:invariants}

For plain 4B (Table~\ref{tab:cj}), \cond{degree} lowers $J$ from 85.75\% to 76.50\%
($-9.25$ [$-13.75$, $-4.75$]; 28 nodes fixed, 65 broken) and $C$ from
87.00\% to 81.00\%. The degree task carries the larger share ($-6.50$,
against $-3.00$ on neighbors). The answer-carrying primer thus buys a leaked
global gain on \task{edge\_count} while local joint correctness falls,
though by no more than under the length control. \cond{filler} lowers $J$
by 13.25 pp, through the neighborhood task ($-12.50$). \cond{components}, a
37-character sentence, raises $J$ to 95.00\% (+9.25 [6.00, 12.75]). For
plain 1.7B, \cond{degree} raises $J$ by 8.50 pp [4.25, 12.75], but $C$ rises
further, to 73.00\% against $J$ of 62.50\%. So 10.50\% of pairs satisfy
$d_v{=}|N(v)|$ with at least one wrong answer: agreement overstates
correctness. For 4B-T, whose $J$ is 93.75\% without a primer, \cond{degree}
adds 2.00 [$-1.25$, 5.00].

\begin{table}[t]
\centering
\small
\begin{tabular}{lrrr}
\toprule
Arm & None & Degree & All \\
\midrule
1.7B   & 62.25/54.00 & 73.00/62.50 & 68.00/61.75 \\
1.7B-T & 81.00/71.00 & 81.00/76.50 & 83.50/79.50 \\
4B     & 87.00/85.75 & 81.00/76.50 & 80.00/77.50 \\
4B-T   & 94.00/93.75 & 96.00/95.75 & 93.75/93.25 \\
\bottomrule
\end{tabular}
\caption{Degree--neighborhood agreement / joint correctness ($C$/$J$, \%;
400 queried nodes per cell, $p\le.50$).}
\label{tab:cj}
\end{table}

\subsection{Side information acts through global cues}
\label{sec:side}

On \task{edge\_existence} the share of queried pairs that are edges rises
with density, from 10\% at $p{=}.10$ to 85\% at $.85$. Pooled over densities,
the plain arms answer ``yes'' to 97--100\% of true edges under every
condition, and 94--99\% of their errors are false alarms. Plain 1.7B's
success without a primer rises from 54\% at $p{=}.50$ to 86\% at $.85$, but
it answers ``yes'' to 94\% and 99\% of queries there. Its balanced success
is 55.77\% and 53.33\%; always answering ``yes'' scores 85\% at $.85$
(Figure~\ref{fig:density}b). Primers move this label bias rather than
edge recall. Pooled over the seven densities, \cond{all} lowers 1.7B's
false-alarm rate from 51\% to 33\% ($-18.6$ [$-24.0$, $-13.1$]).
\cond{filler}, which states nothing about edges, raises it to 69\% (+18.3
[13.9, 22.7]). Where counting degrades, \cond{clustering} adds 11.3 pp [6.7, 16.3] to
plain 4B on \task{node\_degree} above $p{=}.50$. The model keeps enumerating
(69\% of responses, against 74\% without a primer). The printed coefficients
barely vary within a graph (mean within-graph SD at most
0.021), so any gain rests at most on a graph-level value. On these cells
\cond{degree}, \cond{clustering} and \cond{rwse} gain 9.3, 11.3 and 7.7 pp,
but \cond{all}, which combines them, loses 21.3.

\subsection{Two decimals collapse RWSE resolution}
\label{sec:resolution}

In a GNN, RWSE is useful partly because its continuous values can separate
nodes of equal degree, which the first 1-WL round cannot tell apart. In a
prompt, each value becomes a fixed-precision digit string, and values that
round alike become identical token strings. Printed to two decimals, the
$(P^2_{vv}, P^3_{vv})$ pairs
of the 40 nodes take on average 30.14 distinct values at $p{=}.10$, 6.74 at
$.35$, 3.30 at $.50$ and 2.00 at $.85$ (Figure~\ref{fig:density}c). At
$p{=}.50$ the most common pair is shared by 62.68\% of nodes. From
$p{=}.35$ the \cond{degree} primer takes more distinct values than
the RWSE primer (11.98 against 6.74 at $p{=}.35$; 12.50 against 3.30 at
$.50$). There, the printed encoding separates fewer node classes than the
first 1-WL round, although RWSE is meant to add information beyond 1-WL.
The loss occurs at the printed precision, not in the graph: at six decimals the same
features take 40.00 distinct values at $p{=}.50$ and 39.99 at $.85$. Without a higher-precision condition, we
attribute no accuracy change to the loss.

\section{Discussion and Conclusion}

Where a primer carries the answer, a model can change procedure, and the
sign of the effect depends on what is displaced. Where traceable, side
information acts via completion and ``yes'' bias.
Agreement overstates correctness, and a leaked answer can lower joint
correctness while raising a global count. From $p{=}.35$, two-decimal
RWSE separates fewer node classes than 1-WL's first round.

For graph ML evaluation, gains from a structural interface should be reported
with (i)~a graph-blind score, to separate leakage from use; (ii)~per-cell
baselines, on which a leaked answer's effect depends; (iii)~joint correctness and $C{-}J$, not
agreement alone; (iv)~balanced accuracy where density
changes label prevalence; (v)~a content-free length control; and (vi)~the
number of distinct printed feature values per graph. On this corpus,
structural primers trigger procedural shortcuts that help where the model's
own count partly fails, harm where the count succeeds, and act as a check
only for a model that still counts. None of these routes shows general
graph understanding.

\section*{Limitations}

We evaluate two sizes of one model family on synthetic Erd\H{o}s--R\'enyi
graphs, with one encoding, node order, primer position and prompt template.
Node position matters: retrieval is more accurate for nodes 0--9 (96\% against
83\%) \citep{liu2024lostinthemiddle,ge2025graph}. Decoding is greedy. The
extension's 2,048-token plain-arm budget confounds density with budget above
$p{=}.50$, and the thinking arms' completion effects may be specific to our
budgets. Per-density effects rest on 100 graphs each. \cond{filler} is an
approximate, not token-matched, length control. Procedures are classified
by surface patterns, for plain 4B, 1.7B and 1.7B-T only, and a model may
follow a procedure it does not state \citep{turpin2023unfaithful}. A
targeted, non-exhaustive audit of the reused answer extractor (52
responses) moves $C$ and $J$ by at most 1.25 pp per cell. We test
graph-query primitives, not node classification or link prediction.

\section*{Ethics Statement}

Open-weight models on synthetic graphs; no human subjects, personal data or
annotation.

\bibliographystyle{acl_natbib}
\bibliography{custom}

\end{document}
